\documentclass[12pt,a4paper,oneside]{article}
\usepackage[brazil]{babel}
\usepackage[utf8]{inputenc}
\usepackage{olymp}
\usepackage{color}
\usepackage[dvipsnames]{xcolor}
\usepackage{listings}

\setcounter{problem}{2}

\begin{document}

\begin{problem}{Polinômio}{poli}{2}
 
Polinômios são expressões algébricas formadas pela soma algébrica de monômios. Um monômio é expresso na forma $ax^n$, sendo $x$  uma variável, $a$ um número real (o coeficiente) e $n$ um número natural (o grau). O grau de um polinômio é maior grau dos seus monômios.

Por exemplo, um polinômio de grau $3$ e coeficientes $1$ $-2$ $0$ $4$ é representado por:

\begin{equation}
1x^3 - 2x^2 + 0x^1 + 4x^0\notag
\end{equation}

ou seja,

\begin{equation}
x^3 - 2x^2 + 4\notag
\end{equation}

Para $x=1$, este polinômio vale $3$. Para $x=2$, vale $4$. E, para $x=-2$, vale $-12$.

Faça um programa para calcular o valor de um polinômio para um dado valor de variável.

%\Notes
%
%Observações, se tiver.

\InputFile

A entrada é composta por três linhas. A primeira contém um número $N$, o grau do polinômio. A segunda contém $N+1$ valores reais: $A_N, A_{N-1}, \dots, A_1, A_0$, os coeficientes dos monômios do polinômio. A terceira linha contém um valor real $X$, o valor da variável.

Restrições: $2 \leq N \leq 10$, $-100.0 \leq A_i \leq 100.0$, $-10 \leq X \leq 10$.

\OutputFile

Seu programa deve gerar uma linha de saída contendo o valor do polinômio para o valor da variável, formatado para 2 casas decimais.

%\Notes
%Nesta questão, seu programa deve usar a função \texttt{fatorial} dada %acima exatamente como está, não a modifique! Sua
%tarefa é apenas criar o programa com a função \texttt{main}.

\Example

\begin{example}
\exmp{
3
1 -2 0 4
-2
}{
-12.00
}%
\end{example}

\begin{example}
\exmp{
3
1 -2 0 4
0
}{
4.00
}%
\end{example}

\begin{example}
\exmp{
6
2 -4 2 1 -3 4 -5
2
}{
31.00
}%
\end{example}

\begin{example}
\exmp{
6
2 -4 2 1 -3 4 -5
1.5
}{
0.16
}%
\end{example}

\begin{example}
\exmp{
4
1 1.1 1.2 1.3 1.4
4
}{
352.20
}%
\end{example}

\begin{example}
\exmp{
4
1 1.1 1.2 1.3 1.4
-4
}{
201.00
}%
\end{example}

%Comentários sobre o exemplo, se necessário.

\end{problem}

\end{document}
